\chapter{单调算子与预解式}

\label{chap:09}

第 \ref{chap:05} 章把凸泛函的局部最优性写成了次微分条件，第 \ref{chap:06} 章说明共轭与下确界卷积怎样把这些局部信息重新组织成对偶与平滑，第 \ref{chap:07} 章则把原--对偶系统放进锥约束与 KKT 语言中统一处理。到这一章，问题已经不再是“如何写出最优性条件”，而是“如何把最优性条件改写成可迭代求解的对象”。单调算子理论正是完成这一步的标准接口。

本章沿四条线索推进。第 \ref{sec:9-1} 节先在 Banach/Hilbert 的双层背景下定义单调与极大单调，并把次微分、法锥这些熟悉对象提升为极大单调算子的原型。第 \ref{sec:9-2} 节随后把视角收束到 Hilbert 空间：在这里，借助第 \ref{sec:3-1} 节定理~\ref{thm:riesz-representation-hilbert} 的向量化识别，resolvent 可以被具体写成邻近算子，投影、软阈值和 Moreau 分解也都获得清晰的几何意义。第 \ref{sec:9-3} 节继续讨论 Moreau--Yosida 正则化，说明第 \ref{sec:6-3} 节中的 Moreau 包络如何从对偶代数对象变成算法中的平滑代理。最后第 \ref{sec:9-4} 节把最优化、变分不等式和可行性问题统一改写为单调包含
\[
0\in Au,
\]
从而把第 \ref{chap:10} 章的分裂算法全部放到同一个零点框架下。

阅读本章时，最需要记住的是环境分层。对一般 Banach 空间而言，单调算子的自然值域是对偶空间 $X^\ast$；因此“单调”这一概念本身并不依赖 Hilbert 结构。真正把 resolvent、prox、firm nonexpansive 和 Moreau 分解写成最简洁形式的，是 Hilbert 几何。也正因为如此，本章采用双层叙述：第 \ref{sec:9-1} 节保留 Banach/Hilbert 共通的极大单调主线，第 \ref{sec:9-2} 节以后则明确回到 Hilbert 主舞台，不把更一般 Banach 版的近端理论塞进正文中心。

如果把第 \ref{chap:08} 章看成“抽象理论向 Boyd 母语的有限维坍缩”，那么本章恰好走的是相反方向：Boyd 中的 prox、投影、平滑和固定点，在这里被重新解释为极大单调、resolvent、Yosida 正则化和单调包含。这样做的目的不是增加抽象术语，而是为了让第 \ref{chap:10} 章之后的算法不再停留于经验公式，而能被理解为对统一解析结构的计算实现。

\section{极大单调算子}

\label{sec:9-1}

第 \ref{sec:5-3} 节已经告诉我们，不可微凸泛函的局部信息可以写成次微分；第 \ref{sec:7-4} 节则说明，许多最优性系统最终都能写成某种驻点与互补关系。本节要做的，是把这些静态条件提升为动态对象。所谓“动态”，并不是引入时间，而是把“解满足某个不等式”改写成“某个算子取零点”。一旦这一步完成，后面的 resolvent、prox 与分裂算法就都有了统一母体。

\subsection{Banach 空间中的单调与极大单调}

\begin{definition}[单调算子]\label{def:monotone-operator}
设 $X$ 为实 Banach 空间，$A\colon X\rightrightarrows X^\ast$ 为一个多值算子，其图像记为
\[
\operatorname{gra} A:=\{(x,x^\ast)\in X\times X^\ast:x^\ast\in A x\}.
\]
若对任意 $(x,x^\ast),(y,y^\ast)\in \operatorname{gra} A$ 都有
\[
\langle x^\ast-y^\ast,x-y\rangle\ge 0,
\]
则称 $A$ 为\emph{单调算子}。
\end{definition}

\begin{definition}[极大单调算子]\label{def:maximal-monotone-operator}
在定义~\ref{def:monotone-operator} 的条件下，若 $A$ 的图像在集合包含意义下不能被任何更大的单调图像真扩张，即只要某个单调算子 $B$ 满足
\[
\operatorname{gra} A\subset \operatorname{gra} B,
\]
就必有 $\operatorname{gra} A=\operatorname{gra} B$，则称 $A$ 为\emph{极大单调算子}。
\end{definition}

\begin{remark}
单调算子最自然的定义舞台是 $X\times X^\ast$，因此它并不需要 Hilbert 结构。真正需要 Hilbert 几何的是 resolvent、邻近算子和 Moreau 分解，因为这些对象要把“对偶元素”重新识别成“向量”。本章正是按这个分工推进：定义~\ref{def:monotone-operator} 与定义~\ref{def:maximal-monotone-operator} 保持 Banach 级别的一般性，而第 \ref{sec:9-2} 节以后则显式回到 Hilbert 空间。
\end{remark}

\begin{proposition}[单调图像的基本结构]
设 $A\colon X\rightrightarrows X^\ast$ 为单调算子，则对每个 $x\in X$，集合 $A x$ 是凸集。若 $A$ 还是极大单调的，则其图像在积拓扑下闭。
\end{proposition}

\begin{proof}
先证凸值性。若 $x^\ast_1,x^\ast_2\in A x$，取任意 $\theta\in[0,1]$ 并令
\[
x^\ast_\theta:=\theta x^\ast_1+(1-\theta)x^\ast_2.
\]
对任意 $(y,y^\ast)\in \operatorname{gra} A$，由单调性分别有
\[
\langle x^\ast_i-y^\ast,x-y\rangle\ge 0,\qquad i=1,2.
\]
将两式作凸组合即得
\[
\langle x^\ast_\theta-y^\ast,x-y\rangle\ge 0.
\]
若 $A$ 极大单调，则点 $(x,x^\ast_\theta)$ 与整张图像仍保持单调关系。由极大性，必有 $(x,x^\ast_\theta)\in \operatorname{gra} A$，故 $A x$ 为凸集。

再证图像闭性。设 $(x_n,x_n^\ast)\in \operatorname{gra} A$ 且 $(x_n,x_n^\ast)\to (x,x^\ast)$。对任意 $(y,y^\ast)\in \operatorname{gra} A$，单调性给出
\[
\langle x_n^\ast-y^\ast,x_n-y\rangle\ge 0.
\]
令 $n\to\infty$ 得
\[
\langle x^\ast-y^\ast,x-y\rangle\ge 0.
\]
于是 $(x,x^\ast)$ 与图像中的每个点都保持单调关系。由定义~\ref{def:maximal-monotone-operator} 的极大性，$(x,x^\ast)\in \operatorname{gra} A$。
\end{proof}

\subsection{Hilbert 识别与 resolvent}

\begin{definition}[单调算子的 resolvent]\label{def:resolvent-monotone-operator}
设 $H$ 为实 Hilbert 空间，$A\colon H\rightrightarrows H$ 为单调算子，$\lambda>0$。称多值映射
\[
J_{\lambda A}:=(I+\lambda A)^{-1}
\]
为 $A$ 的\emph{resolvent}。也就是说，
\[
u=J_{\lambda A}(x)
\quad\Longleftrightarrow\quad
x-u\in \lambda A u.
\]
\end{definition}

\begin{theorem}[Minty 满射定理]\label{thm:minty-surjectivity}
设 $H$ 为 Hilbert 空间，$A\colon H\rightrightarrows H$ 为单调算子，$\lambda>0$。则在 $\operatorname{ran}(I+\lambda A)$ 上，resolvent $J_{\lambda A}$ 是单值的，并满足参数化公式
\[
\operatorname{gra} A
=
\left\{
\left(J_{\lambda A}(x),\frac{x-J_{\lambda A}(x)}{\lambda}\right):
x\in \operatorname{ran}(I+\lambda A)
\right\}.
\]
进一步，$A$ 为极大单调算子，当且仅当
\[
\operatorname{ran}(I+\lambda A)=H.
\]
\end{theorem}

\begin{proof}
先证单值性。若 $u_1,u_2\in J_{\lambda A}(x)$，则存在 $a_i\in A u_i$ 使
\[
x=u_i+\lambda a_i,\qquad i=1,2.
\]
两式相减得
\[
u_1-u_2=-\lambda(a_1-a_2).
\]
再与 $a_1-a_2$ 作内积并用单调性，
\[
\|u_1-u_2\|^2
=
-\lambda\langle a_1-a_2,u_1-u_2\rangle
\le 0.
\]
故 $u_1=u_2$，于是 $J_{\lambda A}$ 在其定义域上单值。

参数化公式只是定义~\ref{def:resolvent-monotone-operator} 的改写。若 $(u,a)\in \operatorname{gra} A$，取
\[
x:=u+\lambda a,
\]
则 $u=J_{\lambda A}(x)$ 且 $a=(x-u)/\lambda$。反过来，若 $x\in \operatorname{ran}(I+\lambda A)$，令 $u=J_{\lambda A}(x)$，则由 resolvent 的定义，
\[
\frac{x-u}{\lambda}\in A u,
\]
故参数化成立。

最后说明极大性与满射性的关系。若 $\operatorname{ran}(I+\lambda A)=H$，则对每个 $x\in H$，存在唯一 $u=J_{\lambda A}(x)$。若某个点 $(z,w)$ 与 $\operatorname{gra} A$ 的每个点都保持单调关系，即
\[
\langle w-a,z-u\rangle\ge 0,\qquad \forall (u,a)\in \operatorname{gra} A,
\]
取 $x:=z+\lambda w$。由满射性可取 $u=J_{\lambda A}(x)$，并令 $a=(x-u)/\lambda\in A u$。于是
\[
z-u=x-\lambda w-u=\lambda(a-w).
\]
代回单调关系得
\[
0\le \langle w-a,z-u\rangle
=
-\lambda\|w-a\|^2,
\]
故 $w=a$，继而 $z=u$，从而 $(z,w)\in \operatorname{gra} A$，即 $A$ 极大。

反之，若 $A$ 极大单调，则可对图像施加标准的 Minty 变换
\[
(u,a)\mapsto (u+\lambda a,u-\lambda a).
\]
单调性恰好保证该变换后的第一坐标决定第二坐标，而极大性保证其第一坐标像不能遗漏任何点，否则可用 Hahn--Banach 分离构造出对图像的真单调扩张，违背定义~\ref{def:maximal-monotone-operator}。因此
\[
\operatorname{ran}(I+\lambda A)=H.
\]
这就是 Minty 满射定理。
\end{proof}

\subsection{次微分作为极大单调算子的原型}

\begin{theorem}[次微分的极大单调性]\label{thm:subdifferential-maximal-monotone}
设 $H$ 为 Hilbert 空间，$f\colon H\to \mathbb R\cup\{+\infty\}$ 为 proper、凸且下半连续的泛函。则其次微分算子
\[
\partial f\colon H\rightrightarrows H
\]
是极大单调算子。
\end{theorem}

\begin{proof}
由定义~\ref{def:subdifferential}，若 $u^\ast\in \partial f(u)$，$v^\ast\in \partial f(v)$，则
\[
f(v)\ge f(u)+\langle u^\ast,v-u\rangle,
\]
\[
f(u)\ge f(v)+\langle v^\ast,u-v\rangle.
\]
两式相加即得
\[
\langle u^\ast-v^\ast,u-v\rangle\ge 0,
\]
故 $\partial f$ 单调。

为证极大性，由定理~\ref{thm:minty-surjectivity} 只需证明
\[
\operatorname{ran}(I+\partial f)=H.
\]
任取 $y\in H$，考虑泛函
\[
\Phi_y(x):=f(x)+\frac12\|x-y\|^2.
\]
由第 \ref{sec:5-1} 节定义~\ref{def:lower-semicontinuity-functional}，$\Phi_y$ 仍下半连续；二次项提供强制性，因此第 \ref{sec:5-1} 节定理~\ref{thm:weak-lsc-attainment} 保证 $\Phi_y$ 有极小点 $\bar x$。对该极小点，命题~\ref{prop:subgradient-optimality} 给出
\[
0\in \partial \Phi_y(\bar x).
\]
再由第 \ref{sec:5-4} 节定理~\ref{thm:moreau-rockafellar-sum}，
\[
\partial \Phi_y(\bar x)=\partial f(\bar x)+\bar x-y.
\]
于是
\[
y-\bar x\in \partial f(\bar x),
\]
等价于
\[
y\in \bar x+\partial f(\bar x).
\]
这说明每个 $y\in H$ 都属于 $\operatorname{ran}(I+\partial f)$，故由定理~\ref{thm:minty-surjectivity}，$\partial f$ 极大单调。
\end{proof}

\begin{remark}
定理~\ref{thm:subdifferential-maximal-monotone} 是本章与前几章的真正汇合点。第 \ref{sec:5-3} 节的次微分把不可微凸分析写成局部支撑超平面语言，第 \ref{sec:4-3} 节定理~\ref{thm:hahn-banach-geometric} 解释了这些支撑超平面的来源，而第 \ref{sec:6-4} 节定理~\ref{thm:fenchel-rockafellar-duality} 又告诉我们它们怎样进入原--对偶系统。到了这里，同一个对象被重新理解为极大单调算子，从而获得了迭代求解的接口。
\end{remark}

\subsection{典型例子}

\begin{example}[闭真凸下半连续泛函的次微分]
定理~\ref{thm:subdifferential-maximal-monotone} 表明，对任何 proper、凸且下半连续的泛函 $f$，算子 $\partial f$ 都是极大单调的。把这个例子放在这里，是为了让读者看到：本章并不是另起炉灶，而是在第 \ref{sec:5-3} 节对象上增加了“可做 resolvent 和求零点”的新结构。
\end{example}

\begin{example}[闭凸集的法锥]
若 $C\subset H$ 为非空闭凸集，则第 \ref{sec:5-4} 节推论~\ref{cor:normal-cone-indicator} 说明
\[
N_C(x)=\partial \iota_C(x).
\]
因此法锥也是极大单调算子。把这个例子放在这里，是为了说明约束集合并不只贡献可行域，它们还会在后文分裂算法中以算子形式重新出现。
\end{example}

\begin{example}[正半定矩阵诱导的线性单调算子]
在 $\mathbb R^n$ 中，设 $Q\succeq 0$ 为对称正半定矩阵，并定义 $A x:=Qx$。则
\[
\langle Qx-Qy,x-y\rangle=(x-y)^\top Q(x-y)\ge 0,
\]
所以 $A$ 单调。把这个例子放在这里，是为了给 Boyd 读者一个最直接的有限维坍缩对象：矩阵正半定性在本章恰好对应线性算子的单调性。
\end{example}

\subsection*{有限维对照}

Boyd 书里常把一阶最优性条件、投影和 prox 更新直接写成 Euclidean 公式；从本节角度看，这些对象其实都源于极大单调算子。有限维里，由于 $\mathbb R^n$ 与其对偶空间天然等同，且矩阵正半定性可直接检查，单调结构会被压缩成非常熟悉的线性代数语句。无穷维里则必须先经过定义~\ref{def:monotone-operator}、定义~\ref{def:maximal-monotone-operator} 与定理~\ref{thm:minty-surjectivity} 的整理，后面的 resolvent 和 prox 才有稳定落脚点。

\subsection*{习题（待补）}

\begin{exercise}[Minty 参数化占位]
补一题：从定义~\ref{def:resolvent-monotone-operator} 出发，证明单调算子的 resolvent 在其定义域上单值，并写出定理~\ref{thm:minty-surjectivity} 中的图像参数化公式。
\end{exercise}

\begin{exercise}[法锥与极大单调桥接题占位]
补一题：取一个非空闭凸集 $C$，证明 $N_C=\partial \iota_C$，并据此解释为什么约束集合也能被视作极大单调算子的来源。
\end{exercise}


\section{邻近算子}

\label{sec:9-2}

上一节已经把极大单调算子确立为统一母语；本节要做的，是把这门语言翻译成 Hilbert 几何。翻译的关键是：在 Hilbert 空间里，resolvent 不只是一个抽象逆映射，它还等价于“在原目标上加一个二次正则项后求极小值”的操作。邻近算子因此既是解析对象，也是算法模块。第 \ref{chap:10} 章的几乎所有分裂算法，都可以看成在不同位置调用本节定义的 prox。

\subsection{prox 的定义与 resolvent 解释}

\begin{definition}[邻近算子]\label{def:proximal-mapping}
设 $H$ 为 Hilbert 空间，$f\colon H\to \mathbb R\cup\{+\infty\}$ 为 proper、凸且下半连续的泛函，$\lambda>0$。定义
\[
\operatorname{prox}_{\lambda f}(x)
:=
\arg\min_{u\in H}
\left\{
f(u)+\frac{1}{2\lambda}\|u-x\|^2
\right\},
\qquad x\in H.
\]
该映射称为 $f$ 的\emph{邻近算子}。
\end{definition}

\begin{proposition}[prox 是次微分 resolvent]\label{prop:prox-resolvent-subdifferential}
在定义~\ref{def:proximal-mapping} 的条件下，
\[
\operatorname{prox}_{\lambda f}=(I+\lambda \partial f)^{-1}=J_{\lambda\partial f}.
\]
特别地，邻近算子是第 \ref{sec:9-1} 节定义~\ref{def:resolvent-monotone-operator} 的 resolvent 在次微分算子上的具体化。
\end{proposition}

\begin{proof}
固定 $x\in H$，考虑泛函
\[
\Phi_x(u):=f(u)+\frac{1}{2\lambda}\|u-x\|^2.
\]
二次项连续且严格凸，故由定义~\ref{def:proximal-mapping} 可知 $\Phi_x$ 的极小点唯一。设其为 $u$。由命题~\ref{prop:subgradient-optimality}，
\[
0\in \partial \Phi_x(u).
\]
再由第 \ref{sec:5-4} 节定理~\ref{thm:moreau-rockafellar-sum}，
\[
\partial \Phi_x(u)=\partial f(u)+\frac{1}{\lambda}(u-x).
\]
因此
\[
0\in \partial f(u)+\frac{1}{\lambda}(u-x)
\quad\Longleftrightarrow\quad
x-u\in \lambda \partial f(u).
\]
这正等价于
\[
u\in (I+\lambda\partial f)^{-1}(x).
\]
由唯一性，$u=\operatorname{prox}_{\lambda f}(x)$，故结论成立。
\end{proof}

\begin{proposition}[prox 的最优性条件]\label{prop:prox-optimality-condition}
设 $u=\operatorname{prox}_{\lambda f}(x)$。则下列两条等价：
\[
u=\operatorname{prox}_{\lambda f}(x);
\]
\[
0\in \partial f(u)+\frac{1}{\lambda}(u-x).
\]
等价地，
\[
\frac{x-u}{\lambda}\in \partial f(u).
\]
\end{proposition}

\begin{proof}
这只是命题~\ref{prop:prox-resolvent-subdifferential} 证明中的最优性条件重写。由于
\[
u=\operatorname{prox}_{\lambda f}(x)
\quad\Longleftrightarrow\quad
x-u\in \lambda \partial f(u),
\]
将后式两边同时除以 $\lambda$ 即得所述等价关系。
\end{proof}

\subsection{非扩张性与 Moreau 分解}

\begin{proposition}[prox 的 firm nonexpansive 性]\label{prop:prox-firm-nonexpansive}
设 $f\colon H\to \mathbb R\cup\{+\infty\}$ 为 proper、凸且下半连续泛函，$\lambda>0$。则对任意 $x,y\in H$，
\[
\|\operatorname{prox}_{\lambda f}(x)-\operatorname{prox}_{\lambda f}(y)\|^2
\le
\left\langle
\operatorname{prox}_{\lambda f}(x)-\operatorname{prox}_{\lambda f}(y),\,
x-y
\right\rangle.
\]
因此 $\operatorname{prox}_{\lambda f}$ 尤其是 $1$-Lipschitz 的。
\end{proposition}

\begin{proof}
记
\[
u=\operatorname{prox}_{\lambda f}(x),
\qquad
v=\operatorname{prox}_{\lambda f}(y).
\]
由命题~\ref{prop:prox-optimality-condition}，
\[
\frac{x-u}{\lambda}\in \partial f(u),
\qquad
\frac{y-v}{\lambda}\in \partial f(v).
\]
再由第 \ref{sec:9-1} 节定理~\ref{thm:subdifferential-maximal-monotone} 所给出的单调性，
\[
\left\langle
\frac{x-u}{\lambda}-\frac{y-v}{\lambda},
u-v
\right\rangle\ge 0.
\]
乘以 $\lambda$ 得
\[
\langle x-y,u-v\rangle-\|u-v\|^2\ge 0,
\]
即
\[
\|u-v\|^2\le \langle u-v,x-y\rangle.
\]
再由 Cauchy--Schwarz 不等式，
\[
\|u-v\|^2\le \|u-v\|\,\|x-y\|,
\]
从而 $\|u-v\|\le \|x-y\|$。
\end{proof}

\begin{proposition}[Moreau 分解]\label{prop:moreau-decomposition-prox}
设 $f\colon H\to \mathbb R\cup\{+\infty\}$ 为 proper、凸且下半连续泛函，$\lambda>0$。则对任意 $x\in H$，
\[
x
=
\operatorname{prox}_{\lambda f}(x)
+\lambda\operatorname{prox}_{\lambda^{-1}f^\ast}(x/\lambda).
\]
\end{proposition}

\begin{proof}
记
\[
u:=\operatorname{prox}_{\lambda f}(x).
\]
由命题~\ref{prop:prox-optimality-condition}，
\[
\frac{x-u}{\lambda}\in \partial f(u).
\]
由第 \ref{sec:6-1} 节命题~\ref{prop:fenchel-young-equality-subgradient}，这等价于
\[
u\in \partial f^\ast\!\left(\frac{x-u}{\lambda}\right).
\]
将其改写为
\[
0\in \partial f^\ast\!\left(\frac{x-u}{\lambda}\right)+\lambda\left(\frac{x-u}{\lambda}-\frac{x}{\lambda}\right).
\]
这正是命题~\ref{prop:prox-optimality-condition} 对泛函 $\lambda^{-1}f^\ast$ 和点 $x/\lambda$ 的应用，因此
\[
\frac{x-u}{\lambda}
=
\operatorname{prox}_{\lambda^{-1}f^\ast}(x/\lambda).
\]
整理即可得到 Moreau 分解。
\end{proof}

\begin{remark}
命题~\ref{prop:moreau-decomposition-prox} 是第 \ref{sec:6-1} 节 Fenchel 共轭与第 \ref{sec:9-1} 节 resolvent 语言在 Hilbert 空间中的一次精确握手。它说明 prox 从来不只是“解一个带二次项的小问题”，而是原空间与对偶空间之间的一种显式分解。
\end{remark}

\subsection{典型例子}

\begin{example}[闭凸集投影]
取 $f=\iota_C$，其中 $C\subset H$ 为非空闭凸集。则
\[
\operatorname{prox}_{\lambda \iota_C}(x)
=
\arg\min_{u\in C}\frac{1}{2\lambda}\|u-x\|^2
=
P_Cx.
\]
把这个例子放在这里，是为了直接把第 \ref{sec:3-1} 节定理~\ref{thm:orthogonal-projection} 回收到算法主线：投影就是最基本的 prox。
\end{example}

\begin{example}[$\ell^1$ 软阈值]
在 $\mathbb R^n$ 上取
\[
f(x)=\|x\|_1.
\]
则 $\operatorname{prox}_{\lambda f}$ 逐坐标写成软阈值
\[
(\operatorname{prox}_{\lambda\|\cdot\|_1}(x))_i
=
\operatorname{sign}(x_i)\max\{|x_i|-\lambda,0\}.
\]
把这个例子放在这里，是为了把 Boyd 书中最熟悉的非光滑更新直接识别为 resolvent 的 Euclidean 坍缩。
\end{example}

\begin{example}[核范数与奇异值阈值]
在矩阵空间上取
\[
f(X)=\|X\|_\ast.
\]
则 $\operatorname{prox}_{\lambda f}$ 对应于对奇异值做软阈值。把这个例子放在这里，是为了说明 prox 并不只适用于向量稀疏化；在矩阵补全、低秩恢复和 TV 型模型中，它同样是核心更新模块。
\end{example}

\subsection*{有限维对照}

Boyd 常把 prox 解释成“加一个二次项后求极小值”，并配以投影和软阈值图形。本节说明，这种解释完全正确，但它在有限维里之所以显得自然，是因为第 \ref{sec:3-1} 节定理~\ref{thm:riesz-representation-hilbert} 允许我们把对偶元素直接看成向量，于是第 \ref{sec:9-1} 节的 resolvent 便被压缩成一个 Euclidean 映射。真正的结构核心不是坐标公式，而是命题~\ref{prop:prox-resolvent-subdifferential}、命题~\ref{prop:prox-firm-nonexpansive} 与命题~\ref{prop:moreau-decomposition-prox}。

\subsection*{习题（待补）}

\begin{exercise}[投影即 prox 占位]
补一题：对非空闭凸集 $C$，证明 $\operatorname{prox}_{\lambda\iota_C}=P_C$，并说明命题~\ref{prop:prox-firm-nonexpansive} 如何退化为投影的非扩张性。
\end{exercise}

\begin{exercise}[软阈值桥接题占位]
补一题：直接计算 $\operatorname{prox}_{\lambda\|\cdot\|_1}$，并利用命题~\ref{prop:moreau-decomposition-prox} 写出相应的对偶分解。
\end{exercise}


\section{Moreau 包络}

\label{sec:9-3}

第 \ref{sec:6-3} 节已经从共轭与下确界卷积的角度引入了 Moreau 包络，但在那里，重点还是代数结构和对偶交换。本节要把同一个对象重新放回算法语境：一方面，Moreau-Yosida 正则化提供了“把不可微目标平滑化而不改写极小值集合”的手段；另一方面，它又通过 Yosida 近似把极大单调算子变成单值 Lipschitz 映射，为第 \ref{chap:10} 章的连续和离散迭代提供统一接口。

\subsection{Yosida 近似}

\begin{definition}[Yosida 近似]\label{def:yosida-approximation}
设 $H$ 为 Hilbert 空间，$A\colon H\rightrightarrows H$ 为极大单调算子，$\lambda>0$。定义
\[
A_\lambda:=\frac{1}{\lambda}(I-J_{\lambda A}),
\]
称其为 $A$ 的\emph{Yosida 近似}。当 $A=\partial f$ 时，$A_\lambda$ 也称为 $f$ 的 Moreau--Yosida 正则化梯度。
\end{definition}

\begin{proposition}[Moreau 包络的梯度公式]\label{prop:moreau-envelope-gradient}
设 $f\colon H\to \mathbb R\cup\{+\infty\}$ 为 proper、凸且下半连续泛函，$\lambda>0$。则第 \ref{sec:6-3} 节定义~\ref{def:moreau-envelope} 中的 Moreau 包络
\[
e_\lambda f(x)
=
\inf_{u\in H}
\left\{
f(u)+\frac{1}{2\lambda}\|u-x\|^2
\right\}
\]
在 $H$ 上 Fr\'echet 可微，且
\[
\nabla e_\lambda f(x)
=
\frac{1}{\lambda}\bigl(x-\operatorname{prox}_{\lambda f}(x)\bigr)
=(\partial f)_\lambda(x).
\]
\end{proposition}

\begin{proof}
设
\[
p(x):=\operatorname{prox}_{\lambda f}(x).
\]
由命题~\ref{prop:prox-optimality-condition}，
\[
\frac{x-p(x)}{\lambda}\in \partial f(p(x)).
\]
因此，依据第 \ref{sec:6-1} 节 Fenchel--Young 取等条件，
\[
f(p(x))+f^\ast\!\left(\frac{x-p(x)}{\lambda}\right)
=
\left\langle \frac{x-p(x)}{\lambda},p(x)\right\rangle.
\]
另一方面，对任意 $h\in H$，由 $p(x)$ 对 $e_\lambda f(x)$ 的最优性，
\[
e_\lambda f(x+h)
\le
f(p(x))+\frac{1}{2\lambda}\|p(x)-x-h\|^2.
\]
减去
\[
e_\lambda f(x)=f(p(x))+\frac{1}{2\lambda}\|p(x)-x\|^2
\]
得
\[
e_\lambda f(x+h)-e_\lambda f(x)
\le
\left\langle \frac{x-p(x)}{\lambda},h\right\rangle+\frac{1}{2\lambda}\|h\|^2.
\]
同理，以 $p(x+h)$ 作为 $e_\lambda f(x)$ 的竞争点，可得反向估计
\[
e_\lambda f(x+h)-e_\lambda f(x)
\ge
\left\langle \frac{x-p(x+h)}{\lambda},h\right\rangle-\frac{1}{2\lambda}\|h\|^2.
\]
由命题~\ref{prop:prox-firm-nonexpansive}，$p$ 连续，从而
\[
\frac{x-p(x+h)}{\lambda}\to \frac{x-p(x)}{\lambda}\qquad (h\to 0).
\]
夹逼后得到
\[
e_\lambda f(x+h)-e_\lambda f(x)-\left\langle \frac{x-p(x)}{\lambda},h\right\rangle=o(\|h\|),
\]
故 $e_\lambda f$ Fr\'echet 可微，且梯度正是所述表达式。根据定义~\ref{def:yosida-approximation}，该梯度也就是 $(\partial f)_\lambda$。
\end{proof}

\begin{proposition}[Yosida 近似的 Lipschitz 性]\label{prop:yosida-lipschitz}
设 $A\colon H\rightrightarrows H$ 为极大单调算子，$\lambda>0$。则 $A_\lambda$ 是单值映射，并满足
\[
\|A_\lambda x-A_\lambda y\|
\le
\frac{1}{\lambda}\|x-y\|,
\qquad \forall x,y\in H.
\]
\end{proposition}

\begin{proof}
由定理~\ref{thm:minty-surjectivity}，$J_{\lambda A}$ 在整个 $H$ 上单值，因此定义~\ref{def:yosida-approximation} 立即说明 $A_\lambda$ 单值。再由定义，
\[
A_\lambda x-A_\lambda y
=
\frac{1}{\lambda}\bigl((x-y)-(J_{\lambda A}x-J_{\lambda A}y)\bigr).
\]
对极大单调算子，其 resolvent 是 firm nonexpansive；这可由定理~\ref{thm:minty-surjectivity} 与命题~\ref{prop:prox-firm-nonexpansive} 在一般 $A$ 上的同样证明得到。于是
\[
\|J_{\lambda A}x-J_{\lambda A}y\|^2
\le
\langle J_{\lambda A}x-J_{\lambda A}y,x-y\rangle.
\]
因此
\[
\begin{aligned}
\|A_\lambda x-A_\lambda y\|^2
&=
\frac{1}{\lambda^2}
\|(x-y)-(J_{\lambda A}x-J_{\lambda A}y)\|^2 \\
&=
\frac{1}{\lambda^2}
\left(
\|x-y\|^2+\|J_{\lambda A}x-J_{\lambda A}y\|^2
-2\langle x-y,J_{\lambda A}x-J_{\lambda A}y\rangle
\right) \\
&\le
\frac{1}{\lambda^2}
\left(
\|x-y\|^2-\|J_{\lambda A}x-J_{\lambda A}y\|^2
\right)
\le
\frac{1}{\lambda^2}\|x-y\|^2.
\end{aligned}
\]
取平方根即得结论。
\end{proof}

\subsection{Moreau--Yosida 正则化}

\begin{theorem}[Moreau--Yosida 正则化]\label{thm:moreau-yosida-regularization}
设 $f\colon H\to \mathbb R\cup\{+\infty\}$ 为 proper、凸且下半连续泛函，$\lambda>0$。则 $e_\lambda f$ 具有以下性质：
\begin{enumerate}
  \item $e_\lambda f$ 为凸且处处有限的连续函数；
  \item $e_\lambda f$ Fr\'echet 可微，且其梯度由命题~\ref{prop:moreau-envelope-gradient} 给出；
  \item 对任意 $x\in H$，
  \[
  e_\lambda f(x)\le f(x);
  \]
  \item 若 $0<\lambda_1<\lambda_2$，则
  \[
  e_{\lambda_1}f(x)\ge e_{\lambda_2}f(x),\qquad \forall x\in H,
  \]
  因而当 $\lambda\downarrow 0$ 时，$e_\lambda f$ 由下逼近 $f$。
\end{enumerate}
\end{theorem}

\begin{proof}
第一条中的凸性来自第 \ref{sec:6-3} 节命题~\ref{prop:infimal-convolution-convexity}，因为 $e_\lambda f=f\square q_\lambda$，而二次函数 $q_\lambda(x)=\|x\|^2/(2\lambda)$ 连续且强凸。处处有限与连续性来自竞争点 $u=x$ 给出的上界以及命题~\ref{prop:moreau-envelope-gradient} 的可微性。

第二条正是命题~\ref{prop:moreau-envelope-gradient}。

第三条由取竞争点 $u=x$ 立得：
\[
e_\lambda f(x)
\le
f(x)+\frac{1}{2\lambda}\|x-x\|^2
=
f(x).
\]

第四条只需比较惩罚项。若 $0<\lambda_1<\lambda_2$，则对任意 $u\in H$，
\[
f(u)+\frac{1}{2\lambda_1}\|u-x\|^2
\ge
f(u)+\frac{1}{2\lambda_2}\|u-x\|^2.
\]
两边对 $u$ 取下确界即可得到
\[
e_{\lambda_1}f(x)\ge e_{\lambda_2}f(x).
\]
因此随 $\lambda\downarrow 0$，$e_\lambda f(x)$ 单调增加且始终不超过 $f(x)$，这正是 Moreau--Yosida 正则化的逼近方向。
\end{proof}

\begin{corollary}[Moreau 包络保持极小值集合]\label{cor:moreau-envelope-minimizers}
在定理~\ref{thm:moreau-yosida-regularization} 的假设下，对任意 $\lambda>0$，
\[
\operatorname*{argmin} f
=
\operatorname*{argmin} e_\lambda f.
\]
特别地，
\[
x^\star\in \operatorname*{argmin} f
\quad\Longleftrightarrow\quad
\nabla e_\lambda f(x^\star)=0.
\]
\end{corollary}

\begin{proof}
若 $x^\star$ 极小化 $f$，则对任意 $u$ 有 $f(u)\ge f(x^\star)$，故
\[
f(u)+\frac{1}{2\lambda}\|u-x^\star\|^2\ge f(x^\star).
\]
对 $u$ 取下确界得
\[
e_\lambda f(x^\star)\ge f(x^\star).
\]
再结合定理~\ref{thm:moreau-yosida-regularization} 第三条，
\[
e_\lambda f(x^\star)=f(x^\star).
\]
因此 $x^\star$ 极小化 $e_\lambda f$。反之，若 $x^\star$ 极小化 $e_\lambda f$，则对任意 $x$，
\[
e_\lambda f(x^\star)\le e_\lambda f(x)\le f(x).
\]
另一方面，
\[
e_\lambda f(x^\star)\le f(x^\star),
\]
故
\[
f(x^\star)\le f(x),\qquad \forall x\in H.
\]
于是 $x^\star$ 也是 $f$ 的极小点。最后，由命题~\ref{prop:subgradient-optimality} 与可微凸函数的一阶判别，
\[
x^\star\in \operatorname*{argmin} e_\lambda f
\quad\Longleftrightarrow\quad
\nabla e_\lambda f(x^\star)=0.
\]
\end{proof}

\subsection{典型例子}

\begin{example}[指标函数的包络与距离平方]
设 $C\subset H$ 为非空闭凸集，$f=\iota_C$。则
\[
e_\lambda \iota_C(x)
=
\frac{1}{2\lambda}\operatorname{dist}(x,C)^2.
\]
把这个例子放在这里，是为了把投影、可行性误差与平滑代理统一起来：当目标只是“离集合还有多远”时，Moreau 包络直接变成距离平方。
\end{example}

\begin{example}[绝对值与 Huber 平滑]
在 $\mathbb R$ 上取 $f(t)=|t|$，则第 \ref{sec:6-3} 节例子已经说明
\[
e_\lambda f(x)=
\begin{cases}
\dfrac{x^2}{2\lambda}, & |x|\le \lambda,\\[4pt]
|x|-\dfrac{\lambda}{2}, & |x|>\lambda.
\end{cases}
\]
把这个例子放在这里，是为了让读者看到 Boyd 中非常熟悉的 Huber 平滑，其实正是 Moreau--Yosida 正则化的最简单坍缩。
\end{example}

\begin{example}[稳健回归与图像去噪]
在稳健回归或变分图像去噪中，原始正则项往往不可微。把它替换成 Moreau 包络后，目标函数虽然被平滑化，但推论~\ref{cor:moreau-envelope-minimizers} 表明极小值结构没有被粗暴破坏。把这个例子放在这里，是为了说明 Moreau 包络不是“方便计算的临时近似”，而是带有几何控制的结构性代理。
\end{example}

\subsection*{有限维对照}

Boyd 书中关于平滑化、Huber 损失与近端算法的讨论，通常以有限维图像和矩阵模型呈现。本节说明，这些现象背后的统一对象其实是定义~\ref{def:yosida-approximation}、命题~\ref{prop:moreau-envelope-gradient} 与定理~\ref{thm:moreau-yosida-regularization}。有限维里，Yosida 正则化看上去像“把尖角磨平”；无穷维里，它同时也是把多值极大单调算子替换为单值 Lipschitz 映射的解析机制。

\subsection*{习题（待补）}

\begin{exercise}[距离平方占位]
补一题：证明当 $f=\iota_C$ 时，Moreau 包络等于 $\operatorname{dist}(\cdot,C)^2/(2\lambda)$，并解释其梯度与投影之间的关系。
\end{exercise}

\begin{exercise}[Huber 平滑桥接题占位]
补一题：对 $f(t)=|t|$ 直接计算 $e_\lambda f$ 与 $\nabla e_\lambda f$，并说明它如何连接软阈值与平滑优化。
\end{exercise}


\section{算子求根问题}

\label{sec:9-4}

如果说前几节完成的是“把局部最优性语言改写成算子语言”，那么本节要做的就是最后一步：把不同类型的问题统一写成零点问题。之所以这一步重要，是因为一旦把目标表述为
\[
0\in Au,
\]
最优化、变分不等式、可行性和 KKT 系统都能共享同一套迭代母语。第 \ref{chap:10} 章的前向--后向分裂、Douglas--Rachford 和 ADMM，本质上都只是围绕这种零点表达所设计的不同求解策略。

\subsection{单调包含与最优化}

\begin{definition}[单调包含问题]\label{def:monotone-inclusion-problem}
设 $H$ 为 Hilbert 空间，$A\colon H\rightrightarrows H$ 为极大单调算子。称问题
\[
0\in Au
\]
为与 $A$ 相关的\emph{单调包含问题}，其解集记为
\[
\operatorname{zer} A:=\{u\in H:0\in Au\}.
\]
\end{definition}

\begin{proposition}[最优化问题是单调包含的特例]\label{prop:optimization-as-monotone-inclusion}
设 $f\colon H\to \mathbb R\cup\{+\infty\}$ 为 proper、凸且下半连续的泛函。则对任意 $\bar x\in H$，
\[
\bar x\in \operatorname*{argmin} f
\quad\Longleftrightarrow\quad
0\in \partial f(\bar x).
\]
因此，凸最小化问题可以统一改写为单调包含
\[
0\in \partial f(x).
\]
\end{proposition}

\begin{proof}
这正是第 \ref{sec:5-3} 节命题~\ref{prop:subgradient-optimality} 的内容。若 $\bar x$ 极小化 $f$，则
\[
0\in \partial f(\bar x).
\]
反过来，若
\[
0\in \partial f(\bar x),
\]
则次梯度不等式给出
\[
f(y)\ge f(\bar x)+\langle 0,y-\bar x\rangle=f(\bar x),\qquad \forall y\in H,
\]
故 $\bar x$ 为极小点。
\end{proof}

\subsection{零点与不动点}

\begin{proposition}[resolvent 不动点与零点等价]\label{prop:resolvent-fixed-point-equivalence}
设 $A\colon H\rightrightarrows H$ 为极大单调算子，$\lambda>0$。则对任意 $u\in H$，
\[
u\in \operatorname{zer} A
\quad\Longleftrightarrow\quad
u=J_{\lambda A}(u).
\]
\end{proposition}

\begin{proof}
由定义~\ref{def:resolvent-monotone-operator}，
\[
u=J_{\lambda A}(u)
\quad\Longleftrightarrow\quad
u-u\in \lambda A u
\quad\Longleftrightarrow\quad
0\in A u.
\]
这正是 $u\in \operatorname{zer} A$。
\end{proof}

\begin{proposition}[变分不等式也是单调包含]\label{prop:variational-inequality-as-inclusion}
设 $C\subset H$ 为非空闭凸集，$B\colon H\to H$ 为单值算子。则下列两条等价：
\begin{enumerate}
  \item $\bar x\in C$ 且
  \[
  \langle B\bar x,y-\bar x\rangle\ge 0,\qquad \forall y\in C;
  \]
  \item
  \[
  0\in B\bar x+N_C(\bar x).
  \]
\end{enumerate}
\end{proposition}

\begin{proof}
由第 \ref{sec:5-4} 节推论~\ref{cor:normal-cone-indicator}，
\[
\xi\in N_C(\bar x)
\quad\Longleftrightarrow\quad
\langle \xi,y-\bar x\rangle\le 0,\qquad \forall y\in C.
\]
因此
\[
0\in B\bar x+N_C(\bar x)
\]
等价于存在 $\xi\in N_C(\bar x)$ 使 $\xi=-B\bar x$。代回上式，便得到
\[
\langle -B\bar x,y-\bar x\rangle\le 0,\qquad \forall y\in C,
\]
即
\[
\langle B\bar x,y-\bar x\rangle\ge 0,\qquad \forall y\in C.
\]
反向推理同理成立。
\end{proof}

\begin{remark}
命题~\ref{prop:optimization-as-monotone-inclusion} 与命题~\ref{prop:variational-inequality-as-inclusion} 说明，优化与变分不等式在第 \ref{sec:9-1} 节极大单调语言里并不是两类平行对象。前者对应 $\partial f$，后者对应 $B+N_C$。第 \ref{chap:07} 章的 KKT 系统和第 \ref{sec:7-2} 节抽象标准型，到了本节都可以被看成更大的单调包含系统。
\end{remark}

\subsection{proximal point 作为统一模板}

\begin{corollary}[proximal point 模板]\label{cor:proximal-point-template}
设 $A\colon H\rightrightarrows H$ 为极大单调算子，$\lambda_k>0$。则求解单调包含
\[
0\in Au
\]
的一个基本迭代模板是
\[
u_{k+1}=J_{\lambda_k A}(u_k).
\]
若特别地 $A=\partial f$，则由命题~\ref{prop:prox-resolvent-subdifferential}，
\[
u_{k+1}=\operatorname{prox}_{\lambda_k f}(u_k).
\]
\end{corollary}

\begin{proof}
由命题~\ref{prop:resolvent-fixed-point-equivalence}，单调包含的解恰好是 resolvent 的不动点。因此最自然的迭代，就是不断应用 $J_{\lambda_k A}$ 逼近其不动点。若 $A=\partial f$，则第一个式子立刻退化为 prox 迭代。
\end{proof}

\subsection{典型例子}

\begin{example}[极小化问题的零点表达]
对任意 proper、凸且下半连续的泛函 $f$，
\[
\min_x f(x)
\qquad\Longleftrightarrow\qquad
0\in \partial f(x).
\]
把这个例子放在这里，是为了把“求极小值”彻底改写成“求零点”，从而让第 \ref{sec:10-1} 节的前向--后向分裂不再依赖特定目标函数形式。
\end{example}

\begin{example}[线性方程组的固定点坍缩]
若 $M\colon H\to H$ 为可逆线性算子，则解方程
\[
Mx=b
\]
等价于求零点
\[
0\in (Mx-b).
\]
再把该算子写成 $A x:=Mx-b$，便可把线性代数中的固定点迭代理解为单调包含框架的有限维坍缩。把这个例子放在这里，是为了说明“算子求根”并不是新奇对象，而是把熟悉的方程论纳入统一语法。
\end{example}

\begin{example}[约束平衡与 KKT 系统]
对第 \ref{sec:7-4} 节定义~\ref{def:kkt-cone-system} 的广义 KKT 条件，可将原变量与乘子并在一起，构造一个块算子，使整个原--对偶系统写成单调包含。把这个例子放在这里，是为了把第 \ref{sec:7-4} 节推论~\ref{cor:kkt-via-fenchel-subdifferential} 与第 \ref{chap:10} 章的原--对偶算法直接接起来：后文分裂的对象不是“某个优化问题”，而是一个更大的零点系统。
\end{example}

\subsection*{有限维对照}

Boyd 在算法章节里常把固定点、prox 更新和原--对偶步骤分开介绍；从本节角度看，它们都只是“求某个单调算子零点”的不同表面形式。有限维里，单调包含往往被压缩成矩阵等式、投影步骤或互补松弛方程，因此看起来像几套彼此独立的方法。本节的作用，是把这些方法重新挂回同一个接口上，好让第 \ref{sec:10-1} 节和第 \ref{chap:10} 章的所有分裂算法都能被视作同一问题族的不同求解器。

\subsection*{习题（待补）}

\begin{exercise}[优化到零点占位]
补一题：把一个具体凸最小化问题改写为 $0\in \partial f(x)$，并验证命题~\ref{prop:resolvent-fixed-point-equivalence} 在该例中的不动点解释。
\end{exercise}

\begin{exercise}[变分不等式桥接题占位]
补一题：对一个闭凸集上的变分不等式，证明它等价于命题~\ref{prop:variational-inequality-as-inclusion} 给出的单调包含，并说明该表达如何通向第 \ref{sec:10-1} 节的算法接口。
\end{exercise}

